\documentclass[10pt]{article} 
\renewcommand{\thempfootnote}{\arabic{mpfootnote}}

\usepackage[hidelinks]{hyperref}

\usepackage[utf8]{inputenc}

\usepackage[T1]{fontenc}
\usepackage[ttdefault=true]{AnonymousPro}

\usepackage{multicol}
\setlength{\columnsep}{1cm}

\usepackage{abc}
\usepackage{graphicx}
\usepackage{musicography}

\usepackage[margin=0.5in]{geometry} 
\geometry{letterpaper}

\usepackage{fancyhdr} 
\pagestyle{fancy} % options: empty , plain , fancy
\renewcommand{\headrulewidth}{0pt} % customise the layout...
\lhead{}\chead{}\rhead{}
\lfoot{}\cfoot{\thepage}\rfoot{}

\begin{document}

\thispagestyle{empty}

\begin{abc}[name=sawyer,options={-O= +c -F single.fmt}]
X:1
T:Mississippi Sawyer
M:C|
L:1/8
K:D
"D"a2af a2af | a2aa b2a2 | "G"g2ge g2ge | g2gg a2g2 | 
"D"f2fg f2e2 | defg a2a2 | "A"ABcd egfe | "D"d2dd d4 :| 
"D"d2f2 d2f2 | defg a2a2 | "A"A2cB A2cB | ABcd egfe | 
"D"d2f2 d2f2 | defg a2a2 | "A"ABcd egfe | "D"d2dd d4 :|
\end{abc}
\vskip 0.25in
\noindent\makebox[\textwidth][c]{%
  \begin{minipage}{7.0in}
    \setlength{\parindent}{1.5em}

    The popular Old Time tune, Mississippi Sawyer, is thought to have
    evolved from the British Isles tune, The Downfall of Paris, and is
    still known by that name in some parts of America, mainly
    Tennessee and the Ozark
    region.\footnote{\url{https://tunearch.org/wiki/Annotation:Downfall\_of\_Paris\_(The)}}
    The tune apparently goes back further in that it can be traced
    back to the French Revolution (c. 1780) as \textit{Ça Ira} or \textit{les
      aristocrats a la lanterne} (translated as `aristocrats to the
    lamp-post' roughly `let's lynch the aristocrats').\footnote{It is
      not often that tunes are popular among both the Irish and
      English, but when explained through the lens of `violent
      revolutionary' it starts to make sense. You may find it
      interesting to listen to The Dubliner's version of The Downfall
      of Paris and think to yourself, `This sweet loping tune was a
      bloody revolution song?! That’s crazy talk!'} In America, it
    appears in print in Knauff’s tune collection, \textbf{Virginia
      Reels} Vol 4, No 4 (Baltimore, 1839) as Love from the Heart, but
    notably Knauff’s tune Mississippi Sawyer in Vol 1, No 7, is a
    different tune all together.\footnote{Unfortunately, I could not
      find a facsimile of the complete Knauff's \textbf{Virginia
        Reels} available online. But all four volumes are available in
      the appendix of Goertzen's \textbf{George P. Knauff's Virginia
        Reels and the History of American Fiddling} (Jackson, MS,
      2017).}

    
    A `Mississippi Sawyer' is the term for submerged tree stumps that
    were uprooted by flooding along the Mississippi River. These
    stumps then catch on passing riverboats making a difficult journey
    that much more hazardous. American writer and humorist, Mark Twain
    (Samuel Langhorne Clemens, 1835 – 1910), named his famous
    fictional character, Tom Sawyer, after this river
    menace.\footnote{\url{https://tunearch.org/wiki/Annotation:Mississippi\_Sawyer\_(1)}}

---Kitty Steetle
  \end{minipage}
}


\end{document}
